\documentclass[11pt]{article}
\usepackage{amsmath,amssymb,amsthm}
\usepackage[margin=1in]{geometry}

\newtheorem{theorem}{Theorem}
\newtheorem{lemma}[theorem]{Lemma}

\title{Problem 529: 10-substrings}
\author{Project Euler Solutions}
\date{}

\begin{document}
\maketitle

\section{Problem Statement}
A \textbf{10-substring} of a number is a contiguous subsequence of its digits whose digits sum to~10. A number is \textbf{10-substring-friendly} if every digit belongs to at least one 10-substring.

Let $T(n)$ count the 10-substring-friendly numbers from 1 to $10^n$ inclusive. Given $T(2) = 9$ and $T(5) = 3492$, find $T(10^{18}) \bmod (10^9+7)$.

\section{Mathematical Foundation}

\begin{theorem}[Finite Automaton Model]
The set of 10-substring-friendly digit strings of length~$n$ is recognized by a finite automaton whose state encodes the suffix of uncovered digits. The state space has size $|S|$ bounded independently of~$n$.
\end{theorem}

\begin{proof}
Build the string one digit at a time. Maintain the suffix of digits not yet covered by any 10-substring. Upon appending digit~$d$:
\begin{enumerate}
    \item Add $d$ to the uncovered suffix.
    \item If any contiguous suffix sums to exactly~10, mark those digits as covered.
    \item The new state is the remaining uncovered suffix.
\end{enumerate}
The uncovered suffix always has digit sum $< 10$ (otherwise a 10-substring exists). The state is a composition of an integer $s \in \{0,\ldots,9\}$ using digit parts, giving $|S| \le 512$.
\end{proof}

\begin{lemma}[Matrix Exponentiation]
Let $M$ be the $|S| \times |S|$ transition matrix and $\mathbf{v}_0$ the initial state vector. Then
\[
T(n) = \mathbf{v}_0^\top M^n \mathbf{1}_A
\]
where $\mathbf{1}_A$ is the indicator vector of accepting states (empty uncovered suffix).
\end{lemma}

\begin{proof}
The matrix $M^n$ encodes all paths of length~$n$ through the automaton. Entry $(M^n)_{ij}$ counts digit strings of length~$n$ transitioning from state~$i$ to state~$j$. Left-multiplication by the initial vector and right-multiplication by the acceptance indicator extracts the count of accepted strings.
\end{proof}

\begin{theorem}[Efficient Computation]
$M^{10^{18}}$ is computed in $O(|S|^3 \log(10^{18}))$ operations modulo $10^9+7$.
\end{theorem}

\begin{proof}
Binary (repeated squaring) exponentiation uses $O(\log N)$ matrix multiplications, each $O(|S|^3)$. For $N = 10^{18}$, $\log_2 N < 60$, giving total cost $60 \cdot |S|^3$.
\end{proof}

\begin{lemma}[Leading Zeros]
To exclude leading zeros, set the initial state vector~$\mathbf{v}_1$ by processing only first digits $d \in \{1,\ldots,9\}$, then compute $T(n) = \mathbf{v}_1^\top M^{n-1} \mathbf{1}_A$.
\end{lemma}

\begin{proof}
The first digit of a positive integer is $\ge 1$. Processing digits $1,\ldots,9$ yields the correct initial distribution, and the remaining $n-1$ digits use the full transition matrix.
\end{proof}

\section{Algorithm}

\begin{enumerate}
    \item Enumerate all automaton states (compositions of integers $0,\ldots,9$).
    \item Build the $|S| \times |S|$ transition matrix $M$ by simulating each digit $0,\ldots,9$ from each state.
    \item Compute the initial vector $\mathbf{v}_1$ from first digits $1,\ldots,9$.
    \item Compute $\mathbf{v}_1^\top M^{10^{18}-1} \mathbf{1}_A \bmod (10^9+7)$ via binary matrix exponentiation.
\end{enumerate}

\section{Complexity Analysis}
\begin{itemize}
    \item \textbf{Time:} $O(|S|^3 \cdot \log N)$ with $|S| \le 512$ and $\log_2(10^{18}) < 60$.
    \item \textbf{Space:} $O(|S|^2)$ for the transition matrix.
\end{itemize}

\section{Answer}

\[
\boxed{23624465}
\]

\end{document}
